% !TEX root = ../main.tex
\section{Fazit}
\label{sec:fazit}

Die Fragen der Versuchsanleitung lassen sich anhand der Messungen wie folgt beantworten:

\begin{itemize}
  \item \emph{Wie beeinflusst man $P$ und $Q$?} Die Wirkleistung wird über den Winkel $\Th$ eingestellt ($P\propto\sin\Th$), die Blindleistung über die Amplitude $U_{\mathrm{WR1}}$, am Versuch also über $\Udc$ (Abschnitte~\ref{sec:vorbereitung} und \ref{sec:messung612}).
  \item \emph{Welche Harmonischen treten bei Grundfrequenztaktung auf?} Es treten nur ungerade Ordnungen auf. $U_n$ fällt mit $1/n$, $I_n$ mit $1/n^2$ ab; die Messung bestätigt die Theorie (Tabelle~\ref{tab:611}).
  \item \emph{Einfluss von $\Th$ und $\Udc$ auf die Harmonischen?} Theoretisch ist $I_n$ proportional zu $\Udc$ und unabhängig von $\Th$; dies wurde nicht gemessen.
  \item \emph{Stromform und Oberschwingungen je Pulsmuster?} Bei \SI{1.15}{\kilo\watt} und $Q_1\approx0$ hatte «3.\ und 5.\ eliminiert» die kleinsten Stromoberschwingungen ($I_3=\SI{9}{\percent}$, $I_5=\SI{7}{\percent}$); die Trägerverfahren brauchten rund \SI{360}{\volt} statt \SI{282}{\volt} Zwischenkreisspannung (Abschnitt~\ref{sec:messung621}). Die Theorie erwartet für das feine Trägerverfahren die kleineren Oberschwingungen; dieser Widerspruch bleibt wegen der stark schwankenden Anzeige als Beobachtung stehen.
  \item \emph{Funktioniert die $P$/$Q$-Einstellung bei den anderen Pulsmustern gleich?} Der gleiche Winkel genügte für alle Muster am Lastpunkt; $\Udc$ muss höher gewählt werden. Eine vollständige Messreihe fehlt (Abschnitt~\ref{sec:messung622}).
  \item \emph{Optimaler $\cos\varphi$ und $\lambda$?} Nicht gemessen.
  \item \emph{Verläufe an den Halbleitern?} Erwartung als Skizze, Messung als Oszillogramm: komplementäre Rechteckspannungen mit dem Hub $\Udc$, Strom mit Steigungswechsel bei jeder Flanke; die Spurzuordnung beruht auf einer Vermutung (Abschnitt~\ref{sec:messung63}).
  \item \emph{Drehzahlbereich der ASM?} Gemessen etwa \SI{1489}{\per\minute} (Motorbetrieb, \SI{1.15}{\kilo\watt}) bis \SI{1505}{\per\minute} (Generatorbetrieb, \SI{1.15}{\kilo\watt}); laut Typenschild \SI{1500}{\per\minute} im Leerlauf und \SI{1430}{\per\minute} bei Nennlast (Abschnitt~\ref{sec:messung64}).
\end{itemize}
